% =========================================================
% Section: Introduction to Lambda Calculus
% Chapter: Lambda Calculus — Volume I
% =========================================================

% ---------------------------------------------------------
% Part I: Untyped Lambda Calculus
% ---------------------------------------------------------
% =========================================================
% References --- Background Prerequisites
% =========================================================




\subsection*{Untyped Lambda Calculus}
\begin{tcolorbox}[title={Background References}]
\begin{itemize}
  \item \textbf{First-order logic:} Enderton~[1972], Barwise~[1977a].
  \item \textbf{General topology:} Kelly~[1955].
  \item \textbf{Set theory:} Halmos~[1960], van Dalen et al.~[1978].
  \item \textbf{Recursion theory:} Rogers~[1967].
  \item \textbf{Category theory:} Arbib and Manes~[1975], MacLane~[1972].
\end{itemize}
\end{tcolorbox}




% =========================================================
% Definition --- Lambda Terms
% =========================================================

\begin{definitionbox}{Definition --- Lambda Terms}
\begin{definition}[Lambda Terms]
\label{def:lambda-terms}
The set $\Lambda$ of $\lambda$-terms is defined inductively as follows:

\begin{enumerate}
    \item Every variable $x$ is a $\lambda$-term.
    \item If $M$ is a $\lambda$-term and $x$ is a variable, then
    \[
    \lambda x.\,M
    \]
    is a $\lambda$-term.
    \item If $M$ and $N$ are $\lambda$-terms, then
    \[
    MN
    \]
    is a $\lambda$-term.
\end{enumerate}
\end{definition}

\NoLocalDependencies
\end{definitionbox}

\begin{remark*}[Standard quantified statement]
\[
M \in \Lambda
\quad\Longleftrightarrow\quad
M \text{ is generated from variables by finitely many applications of }
\lambda x.(\,\cdot\,) \text{ and } (\,\cdot\,)(\,\cdot\,).
\]
\end{remark*}

\begin{remark*}[Interpretation]
The $\lambda$-terms are the smallest set closed under three constructions:
variables, abstraction (function formation), and application (function
evaluation). This inductive definition plays the same role as recursive
definitions in analysis or algebra: every $\lambda$-term is built from
variables by finitely many applications of abstraction and application.
\end{remark*}

\begin{remark*}[Structural View]
Every $\lambda$-term has a tree structure:
\begin{itemize}
    \item variables are leaves,
    \item abstractions $\lambda x.\,M$ are unary nodes,
    \item applications $MN$ are binary nodes.
\end{itemize}
Thus $\Lambda$ is a freely generated syntactic algebra.
\end{remark*}
